\section{2026-05-19}

\subsection{Existence of Partitions of Unity}

\begin{theorem}
Let $M$ be a smooth manifold, possibly with boundary. Let
\[
  \mathcal X=\{X_\alpha\}_{\alpha\in A}
\]
be an indexed open cover of $M$. Then there exists a smooth partition of unity
subordinate to $\mathcal X$.
\end{theorem}

\begin{proof}
For simplicity, assume first that $M$ has no boundary.

Since each $X_\alpha$ is open in $M$, it is a smooth manifold in its own right. Hence each $X_\alpha$ has a basis of regular coordinate balls. Therefore there exists a basis
\[
  \mathcal B_\alpha
\]
of regular coordinate balls for $X_\alpha$.

Thus
\[
  \bigcup_{\alpha\in A}\mathcal B_\alpha
\]
is a basis for $M$. Since $M$ is a smooth manifold, it is paracompact. Hence the open cover $\mathcal X$ has a countable locally finite refinement
\[
  \{B_i\}_{i=1}^\infty
\]
consisting of elements of
\[
  \bigcup_{\alpha\in A}\mathcal B_\alpha.
\]
We may choose these balls so that the family of closures
\[
  \{\overline{B_i}\}_{i=1}^\infty
\]
is also locally finite.

For each $B_i$, choose an index
\[
  \alpha(i)\in A
\]
such that
\[
  B_i\subset X_{\alpha(i)}.
\]
Choose a smaller coordinate ball
\[
  B_i'\subset B_i
\]
such that
\[
  \overline{B_i'}\subset B_i.
\]
Since $B_i$ is a coordinate ball, choose a smooth coordinate map
\[
  \varphi_i:B_i\to \RR^n
\]
such that
\[
  \varphi_i(B_i)=B_{r_i}(0),
  \qquad
  \varphi_i(B_i')=B_{r_i'}(0),
  \qquad
  r_i'<r_i.
\]

Define a smooth function
\[
  f_i:M\to \RR
\]
by
\[
  f_i(p)=
  \begin{cases}
    H_i(\varphi_i(p)), & p\in B_i,\\
    0, & p\in M\setminus B_i,
  \end{cases}
\]
where
\[
  H_i:\RR^n\to\RR
\]
is a smooth bump function satisfying
\[
  H_i>0
  \quad\text{on }\varphi_i(B_i'),
\]
and
\[
  H_i=0
  \quad\text{on }\RR^n\setminus \varphi_i(B_i).
\]
Then $f_i$ is well-defined and smooth, and
\[
  \operatorname{supp} f_i\subset \overline{B_i}\subset X_{\alpha(i)}.
\]

Now define
\[
  f:M\to\RR
\]
by
\[
  f(p)=\sum_i f_i(p).
\]
Since the family
\[
  \{\overline{B_i}\}_{i=1}^\infty
\]
is locally finite, the sum has only finitely many nonzero terms in a neighborhood of each point. Hence $f$ is smooth.

Moreover, the sets $B_i'$ still cover $M$, so for every $p\in M$ there exists some $i$ such that
\[
  p\in B_i',
\]
and therefore
\[
  f_i(p)>0.
\]
Thus
\[
  f(p)>0
\]
for every $p\in M$.

Define
\[
  g_i(p):=\frac{f_i(p)}{f(p)}.
\]
Then each $g_i$ is well-defined and smooth, and
\[
  0\le g_i(p)\le 1.
\]
Moreover,
\[
  \sum_i g_i(p)
  =
  \frac{\sum_i f_i(p)}{f(p)}
  =
  1.
\]
Also,
\[
  \operatorname{supp} g_i
  =
  \operatorname{supp} f_i
  \subset \overline{B_i}
  \subset X_{\alpha(i)}.
\]

For each $\alpha\in A$, define
\[
  \chi_\alpha:M\to\RR
\]
by
\[
  \chi_\alpha
  :=
  \sum_{\alpha(i)=\alpha} g_i.
\]
Since the family $\{\operatorname{supp}g_i\}$ is locally finite, this sum is locally finite. Hence each $\chi_\alpha$ is smooth.

Furthermore,
\[
  \operatorname{supp}\chi_\alpha
  \subset
  \bigcup_{\alpha(i)=\alpha}\operatorname{supp}g_i
  \subset
  X_\alpha,
\]
and
\[
  \sum_{\alpha\in A}\chi_\alpha
  =
  \sum_i g_i
  =
  1.
\]
Thus
\[
  \{\chi_\alpha\}_{\alpha\in A}
\]
is the desired smooth partition of unity subordinate to $\mathcal X$.
\end{proof}

\subsection{Complex Projective Space}

\begin{example}[Complex projective spaces]
Let
\[
  \CC_*^{n+1}:=\CC^{n+1}\setminus\{0\}.
\]
Define an equivalence relation on $\CC_*^{n+1}$ by
\[
  z\sim \lambda z,
  \qquad
  \lambda\in\CC^*.
\]
The quotient space is called the \vocab{complex projective space}
\[
  \CP^n
  :=
  \CC_*^{n+1}/\sim.
\]
Equivalently,
\[
  \CP^n
  =
  \{[z]\mid z\in \CC_*^{n+1}\}.
\]
For
\[
  z=(z^1,\dots,z^{n+1})\in \CC_*^{n+1},
\]
we write
\[
  [z]=[z^1:z^2:\cdots:z^{n+1}].
\]
Thus
\[
  [z^1:\cdots:z^{n+1}]
  =
  [\lambda z^1:\cdots:\lambda z^{n+1}]
\]
for all
\[
  \lambda\in\CC^*.
\]
\end{example}

For each
\[
  i=1,\dots,n+1,
\]
define
\[
  U_i
  :=
  \{[z]\in\CP^n\mid z^i\neq 0\}.
\]
Then
\[
  \{U_i\}_{i=1}^{n+1}
\]
is an open cover of $\CP^n$.

Define
\[
  \varphi_i:U_i\to\CC^n
\]
by
\[
  \varphi_i([z])
  =
  \left(
    \frac{z^1}{z^i},
    \dots,
    \frac{z^{i-1}}{z^i},
    \frac{z^{i+1}}{z^i},
    \dots,
    \frac{z^{n+1}}{z^i}
  \right).
\]
This is well-defined because multiplying $z$ by a nonzero complex scalar does not change any quotient
\[
  \frac{z^j}{z^i}.
\]

The inverse map is given by
\[
  \varphi_i^{-1}(w^1,\dots,w^n)
  =
  [w^1:\cdots:w^{i-1}:1:w^i:\cdots:w^n].
\]
Hence each
\[
  (U_i,\varphi_i)
\]
is a chart for $\CP^n$.

\begin{remark}
The space $\CP^n$ is naturally a complex manifold of complex dimension $n$.
Equivalently, it is a smooth real manifold of real dimension
\[
  2n.
\]
\end{remark}

\begin{claim*}
The charts
\[
  (U_i,\varphi_i),
  \qquad
  i=1,\dots,n+1,
\]
make $\CP^n$ into a complex manifold.
\end{claim*}

\begin{proof}
We need to check that the transition maps are holomorphic.

For
\[
  i,j\in\{1,\dots,n+1\},
\]
we have
\[
  U_i\cap U_j
  =
  \{[z]\in\CP^n\mid z^i\neq 0,\ z^j\neq 0\}.
\]
The transition map is
\[
  \varphi_j\circ \varphi_i^{-1}
  :
  \varphi_i(U_i\cap U_j)
  \to
  \varphi_j(U_i\cap U_j).
\]

Assume, for simplicity, that
\[
  i<j.
\]
Let
\[
  (w^1,\dots,w^n)\in \varphi_i(U_i\cap U_j).
\]
Under the inverse of the $i$-th chart, this point corresponds to
\[
  [z]
  =
  [w^1:\cdots:w^{i-1}:1:w^i:\cdots:w^n].
\]
In homogeneous coordinates, this means
\[
  z^i=1,
\]
and the other coordinates are the entries of
\[
  (w^1,\dots,w^n).
\]

Since the point lies in
\[
  U_i\cap U_j,
\]
we also have
\[
  z^j\neq 0.
\]
Passing from the $i$-th chart to the $j$-th chart means rescaling the homogeneous coordinates by dividing every coordinate by $z^j$. Thus
\[
  [z^1:\cdots:z^{i-1}:1:z^{i+1}:\cdots:z^{j}:\cdots:z^{n+1}]
\]
is rewritten as
\[
  \left[
    \frac{z^1}{z^j}:
    \cdots:
    \frac{z^{i-1}}{z^j}:
    \frac{1}{z^j}:
    \frac{z^{i+1}}{z^j}:
    \cdots:
    1:
    \cdots:
    \frac{z^{n+1}}{z^j}
  \right].
\]

Therefore the transition map
\[
  \varphi_j\circ\varphi_i^{-1}
\]
is given componentwise by functions of the form
\[
  \frac{z^k}{z^j}
  \qquad\text{or}\qquad
  \frac{1}{z^j}.
\]
Equivalently, in the $i$-th affine coordinates, the components are functions of the form
\[
  \frac{w^a}{w^b}
  \qquad\text{or}\qquad
  \frac{1}{w^b},
\]
on the open set where
\[
  w^b\neq 0.
\]

These functions are holomorphic. Hence
\[
  \varphi_j\circ\varphi_i^{-1}
\]
is holomorphic.

Similarly,
\[
  \varphi_i\circ\varphi_j^{-1}
\]
is holomorphic. Therefore the charts are holomorphically compatible.

Hence $\CP^n$ is a complex manifold.
\end{proof}

\begin{remark}
Since every holomorphic map is smooth as a real map, the same atlas also gives
$\CP^n$ the structure of a smooth real manifold of real dimension
\[
  2n.
\]
\end{remark}

\begin{definition}
The smooth structure on $\CP^n$ determined by the atlas
\[
  \{(U_i,\varphi_i)\}_{i=1}^{n+1}
\]
is called the standard smooth structure on complex projective space.
\end{definition}
